\chapter{Shock-Wave Lithotripsy}

\section*{Introduction \& Clinical Indications}
Shock-wave lithotripsy (SWL), also called extracorporeal shock-wave lithotripsy, uses externally generated focused acoustic energy to fragment selected kidney or ureteric stones so that pieces can pass in urine. It is non-incisional and may be performed as a day procedure, but one or more sessions may be necessary. Selection depends on stone size, location, density, skin-to-stone distance, obstruction, infection, anatomy, bleeding risk, and patient preference.

SWL is not appropriate for an infected obstructed system without prior drainage and treatment. It may be unsuitable with an untreated bleeding disorder, pregnancy, inability to target the stone, certain aneurysms near the treatment path, or an obstruction that prevents fragment passage. Ureteroscopy and PCNL are alternatives.

\section*{Relevant Surgical Anatomy}
The stone is targeted using fluoroscopy or ultrasound. The kidney, ureter, ribs, lung, bowel, spleen, liver, and surrounding tissues must be considered. Lower-pole anatomy may limit fragment clearance, while a narrow or obstructed ureter may cause steinstrasse, a column of fragments that blocks drainage.

\section*{Preoperative Workup \& Preparation}
Imaging confirms stone size and location. Urinalysis, culture, renal function, blood count, pregnancy testing when relevant, and bleeding-risk assessment are selected clinically. Anticoagulants and antiplatelet drugs are managed according to bleeding and thrombosis risk. Consent covers pain, bruising, hematuria, infection, steinstrasse, renal hematoma, incomplete clearance, repeat treatment, and conversion to ureteroscopy or PCNL.

\section*{Surgical Technique \& Procedure}
The patient is positioned so the stone lies in the lithotripter's focal zone. Imaging guides repeated shock waves, with energy and rate adjusted to stone and patient factors. Analgesia or anesthesia is provided as needed. Fragments are allowed to pass, and a stent or preliminary drainage is used selectively rather than routinely.

\section*{Complications \& Risks}
Expected effects include bruising, transient hematuria, flank discomfort, and passage of fragments. Complications include urinary infection or sepsis, ureteric obstruction, steinstrasse, renal or perirenal hematoma, hypertension or renal injury, incomplete fragmentation, and need for further procedures. Prompt assessment is required for fever, rigors, severe pain, inability to pass urine, or persistent heavy bleeding.

\section*{Postoperative Care \& Recovery}
Hydration and analgesia are individualized. Patients may be asked to strain urine for fragments and receive follow-up imaging after an appropriate interval. Stone analysis and metabolic evaluation are useful for recurrent or high-risk stone formers. Stents, when used, require a documented removal plan.

\section*{Outcomes \& Prognosis}
Clearance varies with stone size, density, location, anatomy, and number of sessions. SWL has the advantage of lower invasiveness but may be less effective for dense, large, lower-pole, or poorly targeted stones. The preferred method should be discussed rather than described as universally first-line.

\section*{References}
\begin{enumerate}
  \item European Association of Urology. Guidelines on Urolithiasis.
  \item National Health Service. Kidney Stones: Treatment.
  \item American Urological Association. Surgical Management of Stones: Guideline.
\end{enumerate}

\clearpage
